% Source: https://artofproblemsolving.com/community/c3h2128707_a_simple_geometry_problem
% Stable identity: aops-topic-2128707-q01
The triangles have sides $23,24,25$ and $31,32,33$. In a triangle, the altitude to side $\ell$ is $\dfrac{2K}{\ell}$, where $K$ is the area, so the largest altitude is the one to the shortest side.

\begin{asy}
usepackage("amsmath");
size(7cm);
defaultpen(fontsize(10pt));
pair B1=(0,0), C1=(23,0); real x1=(625-576+529)/46; pair A1=(x1,sqrt(625-x1^2));
pair B2=(30,0), C2=(61,0); real x2=(1089-1024+961)/62; pair A2=(30+x2,sqrt(1089-x2^2));
draw(A1--B1--C1--cycle,linewidth(0.9)); draw(A2--B2--C2--cycle,linewidth(0.9));
pair F1=(A1.x,0), F2=(A2.x,0);
draw(A1--F1,red+dashed+linewidth(0.8)); draw(A2--F2,red+dashed+linewidth(0.8));
draw(F1+(0,1.2)--F1+(1.2,1.2)--F1+(1.2,0),red); draw(F2+(0,1.2)--F2+(1.2,1.2)--F2+(1.2,0),red);
label("$23$",(B1+C1)/2,S); label("$25$",(A1+B1)/2,NW); label("$24$",(A1+C1)/2,NE);
label("$31$",(B2+C2)/2,S); label("$33$",(A2+B2)/2,NW); label("$32$",(A2+C2)/2,NE);
label("$h_1$",(A1+F1)/2,W,red); label("$h_2$",(A2+F2)/2,W,red);
label("\mbox{}",(min(currentpicture,user=true)+max(currentpicture,user=true))/2,black); // keeps the text after the figure black when LaTeX typesets the labels inline
\end{asy}

\textbf{First triangle.} The semiperimeter is $36$, so by Heron's formula
\[K_1=\sqrt{36\cdot13\cdot12\cdot11}=\sqrt{144\cdot429}=12\sqrt{429},\]
and $429=3\cdot11\cdot13$ is squarefree. The largest altitude is $h_1=\dfrac{2K_1}{23}=\dfrac{24}{23}\sqrt{429}$, so $y=24+23+429=476$.

\textbf{Second triangle.} The semiperimeter is $48$, so
\[K_2=\sqrt{48\cdot17\cdot16\cdot15}=\sqrt{2304\cdot85}=48\sqrt{85},\]
and $85=5\cdot17$ is squarefree. The largest altitude is $h_2=\dfrac{2K_2}{31}=\dfrac{96}{31}\sqrt{85}$, so $z=96+31+85=212$.

\textbf{Finally,} $y+z=688$, and $6!-688=5!+h$ gives $h=720-688-120=\boxed{-88}$.
